\section*{Capítulo \ref{ch:b2:biogeochemical-cycles} --- Ciclos biogeoquímicos}

\begin{solution}{exo:b2:biogeochemical-cycles:1}
Um \omterm{def:b2:biogeochemical-cycles:reservoir}{reservatório} é um estoque do elemento (uma massa); um \omterm{def:b2:biogeochemical-cycles:reservoir}{fluxo} é uma
taxa de transferência entre \omterm{def:b2:biogeochemical-cycles:reservoir}{reservatórios} (massa por ano); o \omterm{def:b2:biogeochemical-cycles:reservoir}{tempo de
residência} é a massa do \omterm{def:b2:biogeochemical-cycles:reservoir}{reservatório} dividida por sua saída total, o
tempo médio que um átomo passa nele. Plantas terrestres: $450/120 \approx 4$ anos (mais ou
menos o carbono de uma folha; a \omterm{def:b2:plant-meristems:wood}{madeira} fica por décadas, e a média
esconde a dispersão). Oceano profundo: $37000/100 \approx 400$ anos ---
o tempo antes que uma molécula que afundou volte a ver a superfície.
\end{solution}

\begin{solution}{exo:b2:biogeochemical-cycles:2}
$2.5\times 2.13 = \qty{5.3}{GtC}$ por ano. $10/2.13 = \qty{4.7}{ppm}$
--- se tudo ficasse no ar; com a \omterm{prop:b2:biogeochemical-cycles:carbon}{fração atmosférica} de \qty{45}{\%},
\qty{2.1}{ppm}.
\end{solution}

\begin{solution}{exo:b2:biogeochemical-cycles:3}
Abertura: a fixação biológica pela nitrogenase (cianobactérias,
rizóbios, \emph{Azotobacter}, \emph{Frankia}) e os raios. Retorno:
a \omterm{def:b2:microbial-metabolism:anaerobic}{desnitrificação} (anaeróbios facultativos como \emph{Pseudomonas}
e \emph{Paracoccus} em solos e sedimentos anóxicos) e o anammox
(planctomicetos). Entre os dois: a amonificação do nitrogênio orgânico
em amônio pelos decompositores (fungos, bactérias); a nitrificação do
amônio em nitrito (\emph{Nitrosomonas}, arqueias oxidantes de amônia)
e do nitrito em nitrato (\emph{Nitrobacter}); e a
assimilação do amônio ou do nitrato pelas plantas e pelos micróbios.
\end{solution}

\begin{solution}{exo:b2:biogeochemical-cycles:4}
Em escala geológica, o nitrogênio sempre pode ser obtido do inesgotável
$\mathrm{N_2}$ do ar pelos fixadores, que são favorecidos sempre que
falta nitrogênio; o fósforo não tem tal \omterm{def:b2:biogeochemical-cycles:reservoir}{reservatório}, e seu suprimento é
fixado pelo intemperismo, de modo que o teto de longo prazo da
produtividade é o fósforo. Numa estação, a fixação é lenta e custosa, e
a maioria das plantas não consegue realizá-la, de modo que o nitrogênio
disponível é o que limita o crescimento aqui e agora --- e é por isso
que os fertilizantes são sobretudo nitrogenados.
\end{solution}

\begin{solution}{exo:b2:biogeochemical-cycles:5}
$\tau = 1/0.2 = 5$ anos; $M^{*} = F\tau = \qty{10}{t}$, isto é,
$\qty{1}{g/m^3} = \qty{1000}{mg/m^3}$, fortemente eutrófico. Tempo para
\qty{90}{\%}: $\tau\ln 10 = 11.5$ anos. Reduzir a entrada à metade
reduz à metade o \omterm{def:b2:biogeochemical-cycles:reservoir}{estado estacionário}, para \qty{500}{mg/m^3}, atingido na
mesma escala de tempo.
\end{solution}

\begin{solution}{exo:b2:biogeochemical-cycles:6}
Anos 1960: \qty{0.9}{ppm/ano}, \qty{1.9}{GtC/ano}. Anos 2010: \qty{2.3}{ppm/ano},
\qty{5.0}{GtC/ano}. A taxa de aumento cresceu duas vezes e meia,
acompanhando as emissões.
\end{solution}

\begin{solution}{exo:b2:biogeochemical-cycles:7}
\qty{6}{ppm} equivalem a \qty{13}{GtC}, contra uma produção primária
líquida de \qty{35}{GtC}: o dente de serra é um terço dela. Ele é menor
porque a respiração e a decomposição continuam durante o verão (o dente
de serra registra a absorção líquida, e não a bruta), porque as estações
do hemisfério sul são opostas e em parte se compensam, e porque o oceano
amortece a oscilação.
\end{solution}

\begin{solution}{exo:b2:biogeochemical-cycles:8}
$\qty{200}{kg}/(\qty{28}{g/mol}) = 7100$~mol de $\mathrm{N_2}$; $16\times
7100 = \num{114000}$~mol de ATP; $/30 = 3800$~mol de glicose, \qty{690}{kg}.
Uma cultura que constrói \qty{10}{t} de matéria seca fixa cerca de
\qtyrange{15}{20}{t} de açúcar quando se conta a respiração, de modo que
a fixação lhe custa cerca de \qty{4}{\%} de seus fotoassimilados --- o
preço da independência em relação ao nitrogênio do solo.
\end{solution}

\begin{solution}{exo:b2:biogeochemical-cycles:9}
Razão disponível $\mathrm{N:P} = 32/2.2 = 14.5$, abaixo dos 16 de
Redfield: o nitrato se esgota primeiro, restando \qty{0.2}{\micro mol/kg}
de fosfato. Carbono fixado: $32\times 106/16 =
\qty{212}{\micro mol/kg}$, cerca de \qty{2.5}{mg} de carbono por
quilograma de água.
\end{solution}

\begin{solution}{exo:b2:biogeochemical-cycles:10}
$E^{*} = Q\tau = \qty{500}{GtC}$, \qty{235}{ppm} acima dos 285
pré-industriais: cerca de \qty{520}{ppm}. Pior na realidade porque os
sumidouros não são um único compartimento rápido: o oceano superficial
se satura (sua capacidade de absorver $\mathrm{CO_2}$ cai à medida que seu
carbonato é consumido), o sumidouro terrestre enfraquece com o
aquecimento, e a remoção real tem um componente lento --- a mistura com o
oceano profundo e o intemperismo das rochas --- com constantes de tempo
de mil a cem mil anos, de modo que uma fração do excesso não relaxa em
nenhum $\tau$ de décadas.
\end{solution}

\begin{solution}{exo:b2:biogeochemical-cycles:11}
O que foi medido é a troca do $\mathrm{CO_2}$ atmosférico com o
oceano e a biosfera, e não o decaimento: o $\mathrm{{}^{14}C}$ das bombas foi
diluído nos compartimentos muito maiores do oceano superficial, das
plantas e dos solos, que depois devolveram carbono comum. A meia-vida de
11 anos do excesso (uma constante de tempo exponencial de cerca de 16 anos) é
a escala de tempo dessa troca --- a renovação do carbono da atmosfera em
relação ao oceano superficial e à terra, mais longa que o \omterm{def:b2:biogeochemical-cycles:reservoir}{tempo de
residência} bruto de quatro anos porque boa parte do que sai volta
imediatamente.
\end{solution}

\begin{solution}{exo:b2:biogeochemical-cycles:12}
Nitrogênio: fixação natural de cerca de 200~Tg por ano; fixação humana
(Haber--Bosch 120, culturas 60, combustão 30) de cerca de 210: o \omterm{def:b2:biogeochemical-cycles:reservoir}{fluxo}
dobrou, e os \omterm{def:b2:biogeochemical-cycles:reservoir}{reservatórios} afetados --- solos, rios, mares costeiros, o
$\mathrm{N_2O}$ do ar --- são pequenos e rápidos, de modo que a mudança
aparece em décadas, enquanto o \omterm{def:b2:biogeochemical-cycles:reservoir}{reservatório} de $\mathrm{N_2}$ fica intacto.
Carbono: as emissões humanas de 11~GtC por ano são apenas \qty{5}{\%} da
troca natural bruta de cerca de 200, de modo que o \omterm{def:b2:biogeochemical-cycles:reservoir}{fluxo} recebe apenas um
empurrãozinho; mas os \omterm{def:b2:biogeochemical-cycles:reservoir}{fluxos} naturais se anulam e o humano não, e ele
elevou o \omterm{def:b2:biogeochemical-cycles:reservoir}{reservatório} atmosférico em \qty{50}{\%} --- o pequeno empurrão
tem um grande efeito acumulado sobre um \omterm{def:b2:biogeochemical-cycles:reservoir}{reservatório} com um sumidouro
final lento. Qual ciclo é ``mais perturbado'' depende de se contarem os
\omterm{def:b2:biogeochemical-cycles:reservoir}{fluxos} ou a acumulação.
\end{solution}

\begin{solution}{pb:b2:biogeochemical-cycles:1}
\textbf{1.} $285\times 2.13 = \qty{607}{GtC}$; $412\times 2.13 =
\qty{878}{GtC}$; aumento de \qty{271}{GtC}.
\textbf{2.} $271/700 = 0.39$ (a fração nas últimas décadas, contando as
emissões do uso da terra, fica mais perto de $0.45$).
\textbf{3.} O oceano, cerca de um quarto das emissões, e a terra
(rebrota das florestas e fertilização por $\mathrm{CO_2}$), cerca de um terço;
é o restante dos \qty{430}{GtC} que não estão no ar.
\textbf{4.} $175/38000 = \qty{0.5}{\%}$ --- desprezível para o oceano
inteiro; mas a absorção entrou até agora sobretudo na camada superficial
(\qty{900}{GtC}), onde representa um aumento de um décimo ou mais, e a
química dos carbonatos transforma um pequeno aumento do $\mathrm{CO_2}$ dissolvido num
aumento maior da concentração de íons hidrogênio: o pH superficial caiu
$0.1$, um aumento de \qty{30}{\%} da acidez.
\textbf{5.} Com $Q = Q_0 e^{t/T}$, tente $E = Ae^{t/T}$: $A/T = Q_0 -
A/\tau$, logo $A = Q_0\tau T/(\tau + T)$ e a \omterm{prop:b2:biogeochemical-cycles:carbon}{fração atmosférica}
$\dot E/Q = \tau/(\tau + T)$ é constante. Emissões que crescem
\qty{2}{\%} por ano ($T = 50$) com $\tau = 40$ dão $0.44$: a fração é
constante porque o sumidouro e a fonte crescem juntos.
\textbf{6.} Com $Q$ constante, $E \to Q\tau$ e $\dot E \to 0$: a \omterm{prop:b2:biogeochemical-cycles:carbon}{fração
atmosférica} cai para zero à medida que os sumidouros alcançam a fonte ---
no modelo de uma caixa; na realidade, os sumidouros se saturam e a fração
cai mais lentamente.
\textbf{7.} \qty{271}{GtC}.
\textbf{8.} $E^{*} = 10\times 100 = \qty{1000}{GtC}$: $285 +
1000/2.13 = \qty{755}{ppm}$.
\textbf{9.} $E(t) = 1000 + (271 - 1000)e^{-t/100}$; em $t = 80$,
$E = 1000 - 729\times 0.449 = \qty{672}{GtC}$: \qty{600}{ppm}.
\textbf{10.} $Q = 10(1 - t/50)$: solução particular $E_p = at + b$
com $a = -20$, $b = 3000$ (de $a = 10 - b/\tau$ e $0 = -0.2 -
a/\tau$); $E = 3000 - 20t + (271 - 3000)e^{-t/100}$. Em $t = 50$:
$2000 - 2729\times 0.607 = \qty{344}{GtC}$, \qty{447}{ppm}. Depois $E$
decai livremente: em $t = 80$, $344\,e^{-0.3} = \qty{255}{GtC}$,
\qty{405}{ppm}.
\textbf{11.} \qty{600}{ppm} contra \qty{405}{ppm}: o segundo cenário
devolve o ar, em 2100, quase ao ponto em que está hoje.
\textbf{12.} O $\tau$ único faz os sumidouros continuarem trabalhando no
mesmo ritmo sobre um excesso que diminui; na realidade, os sumidouros
rápidos (a camada de mistura, as florestas em crescimento) já absorveram
o que podiam, e a remoção restante se faz pela lenta mistura com o oceano
profundo e pelo intemperismo, com $\tau$ de séculos a milênios, de modo
que o declínio após a neutralidade líquida é muito mais lento --- a
concentração fica mais ou menos constante durante séculos, em vez de
cair.
\textbf{13.} Cultura, \qty{90}{kg}; desnitrificado, \qty{37.5}{kg};
lixiviado, \qty{22.5}{kg} de nitrogênio por hectare por ano.
\textbf{14.} $\qty{22.5}{kg}/\qty{14}{g/mol} = 1600$~mol de N; carbono
$1600\times 106/16 = \num{10600}$~mol, \qty{128}{kg}.
\textbf{15.} \num{10600}~mol de $\mathrm{O_2}$, \qty{340}{kg}.
\textbf{16.} $\num{10600}/0.25 = \num{42000}$~\unit{m^3}: a lixiviação
de um hectare pode retirar o oxigênio de quarenta mil metros cúbicos de
água de fundo a cada ano --- e o Mississippi drena cem milhões de
hectares.
\textbf{17.} $37.5\times 0.01 = \qty{0.375}{kg}$ de nitrogênio na forma de
$\mathrm{N_2O}$, isto é, $0.375\times 44/28 = \qty{0.59}{kg}$ de $\mathrm{N_2O}$;
$\times 270 = \qty{160}{kg}$ de equivalente $\mathrm{CO_2}$.
\textbf{18.} $150\times 4 = \qty{600}{kg}$ de $\mathrm{CO_2}$: a fabricação
supera o $\mathrm{N_2O}$ do campo na razão de quatro para um, e os dois
juntos somam três quartos de tonelada por hectare por ano.
\textbf{19.} $\tau = 10$ anos; $M^{*} = 5\times 10 = \qty{50}{t}$;
concentração $50/(5\times 10^{7}) = \qty{1}{g/m^3} =
\qty{1000}{mg/m^3}$.
\textbf{20.} Trinta vezes acima do limiar: severamente eutrófico.
\textbf{21.} $1/\tau = 0.1 + 0.2 = 0.3$: $\tau = 3.3$ anos; $M^{*} =
5/0.3 = \qty{16.7}{t}$, \qty{330}{mg/m^3}.
\textbf{22.} $M^{*} = 1/0.3 = \qty{3.3}{t}$, \qty{67}{mg/m^3} --- ainda
eutrófico; a menos de \qty{10}{\%} após $\tau\ln 10 = 7.7$ anos.
\textbf{23.} O sedimento não é um sumidouro de sentido único: o fósforo
armazenado nele durante os anos eutróficos é liberado de volta sob as
condições anóxicas que a própria eutrofização cria (o fosfato ligado ao
ferro se dissolve quando o ferro é reduzido), de modo que o lago se
alimenta durante anos depois de cortada a entrada externa --- a carga
interna, um segundo \omterm{def:b2:biogeochemical-cycles:reservoir}{reservatório} com um longo \omterm{def:b2:biogeochemical-cycles:reservoir}{tempo de residência}.
\textbf{24.} $\qty{5000}{kg}/\qty{31}{g/mol} = 1.6\times 10^{5}$~mol de P;
carbono $\times 106 = 1.7\times 10^{7}$~mol, \qty{205}{t} de carbono por
ano --- \qty{20}{g/m^2} sobre a superfície do lago, se ele tiver
\qty{5}{m} de profundidade, uma proliferação intensa.
\textbf{25.} \omterm{prop:b2:biogeochemical-cycles:carbon}{Fração atmosférica} de $0.39$ ao longo de 1850--2020; em 2100,
\qty{600}{ppm} com as emissões mantidas em \qty{10}{GtC/ano} e
\qty{405}{ppm} com as emissões zeradas até 2070 (um modelo otimista de
$\tau$ único); nitrogênio lixiviado de \qty{22.5}{kg} por hectare
por ano.
\end{solution}
